% Figura redibujada a partir de "Imagenes/Circuito equivalente.PNG" (Tema 5).
% Coordenadas en píxeles del PNG original (433 x 312 px), escaladas para que
% ocupe lo mismo que el PNG en el documento: 169.5 pt x 122.1 pt.
\documentclass[12pt,border=0pt]{standalone}
\usepackage[spanish,es-tabla]{babel}
\usepackage[utf8]{inputenc}
\usepackage{amsmath}
\usepackage{xcolor}
\usepackage{tikz}
\usepackage{pgfplots}
\usepackage{circuitikz}
\usetikzlibrary{babel, circuits.ee}

\definecolor{applegreen}{rgb}{0.55, 0.71, 0.0}
\definecolor{airforceblue}{rgb}{0.36, 0.54, 0.66}
\definecolor{amber}{rgb}{1.0, 0.49, 0.0}

\newcommand{\diff}[3][1]{%
	\IfBlankTF{#1}{\frac{\mathrm{d}{#2}}{\mathrm{d}{#3}}}%
	{\frac{\mathrm{d}^{#1} {#2}}{\mathrm{d}{#3}^{#1}}}%
}%
\newcommand{\pdiff}[3][1]{%
	\IfBlankTF{#1}{\frac{\partial {#2}}{\partial {#3}}}%
	{\frac{\partial^{#1} {#2}}{\partial {#3}^{#1}}}%
}%

\begin{document}
\begin{tikzpicture}[x=0.3915pt, y=-0.3915pt, font=\footnotesize,
    abierto/.style={circle, draw, fill=white, inner sep=0pt, minimum size=4pt},
    nodo/.style={circle, fill, inner sep=0pt, minimum size=4pt},
    fuente/.style={circle, draw, thick, fill=white, inner sep=0pt, minimum size=21pt}]
    \useasboundingbox (0,0) rectangle (433,312);
    % Puerta (en abierto)
    \draw (57,46) -- (122,46);
    \node[abierto] at (57,46) {};
    \node[abierto] at (122,46) {};
    % Nodo de drenador y nodo de fuente
    \draw (172,46) -- (357,46) node[abierto]{};
    \draw (62,276) node[nodo]{} -- (382,276) node[nodo]{};
    % Fuentes de corriente dependientes
    \foreach \x in {172,260} {
        \draw (\x,46) -- (\x,276);
        \node[fuente] at (\x,160) {};
        \draw[-latex] (\x,147) -- (\x,174);
    }
    % Resistencia 1/g_d
    \draw (322,46) -- (322,124) (322,196) -- (322,276);
    \draw (311,124) rectangle (333,196);
    % Corriente de drenador
    \draw[-latex] (350,46) -- (327,46);
    % Etiquetas
    \node at (22,48) {G};
    \node at (383,48) {D};
    \node at (250,300) {S};
    \node at (343,18) {$i_\mathrm{d}$};
    \node at (25,160) {$v_\mathrm{gs}$};
    \node at (405,160) {$v_\mathrm{ds}$};
    \node[left] at (147,175) {$g_\mathrm{m} v_\mathrm{gs}$};
    \node at (222,207) {$g_\mathrm{mb} v_\mathrm{bs}$};
    \node at (365,212) {$1/g_\mathrm{d}$};
\end{tikzpicture}
\end{document}
